% ============================================================
% NAC: Nesterov Accelerated Counterattack
% SIVP (Signal, Image and Video Processing) submission version
% Springer Nature sn-jnl template — December 2024
% ============================================================
\documentclass[pdflatex,sn-basic,Numbered]{sn-jnl}

% --- packages ---
\usepackage{graphicx}
\usepackage{multirow}
\usepackage{amsmath,amssymb,amsfonts}
\usepackage{amsthm}
\usepackage{mathrsfs}
\usepackage[title]{appendix}
\usepackage{xcolor}
\usepackage{textcomp}
\usepackage{booktabs}
\usepackage{algorithm}
\usepackage{algorithmicx}
\usepackage{algpseudocode}
\usepackage{listings}
\usepackage{enumitem}

% --- theorem styles ---
\theoremstyle{thmstyleone}
\newtheorem{theorem}{Theorem}
\newtheorem{proposition}[theorem]{Proposition}

\theoremstyle{thmstyletwo}
\newtheorem{example}{Example}
\newtheorem{remark}{Remark}

\theoremstyle{thmstylethree}
\newtheorem{definition}{Definition}

\raggedbottom

\begin{document}

\title[Nesterov Accelerated Counterattack]{Nesterov Accelerated Counterattack for Test-Time Adversarial Defense of Vision-Language Models}

%% ---- authors ----
\author*[1,2]{\fnm{Zhihao} \sur{Zhang}}\email{zhangzhihao126@outlook.com}

\author[1,2]{\fnm{Yazhi} \sur{Liu}}\email{liuyazhi@ncst.edu.cn}

\affil*[1]{\orgdiv{College of Artificial Intelligence}, \orgname{North China University of Science and Technology}, \orgaddress{\city{Tangshan}, \postcode{063210}, \country{China}}}

\affil[2]{\orgdiv{Hebei Key Laboratory of Industrial Intelligent Perception}, \orgname{North China University of Science and Technology}, \orgaddress{\city{Tangshan}, \postcode{063210}, \country{China}}}

%% ---- abstract ----
\abstract{
Vision-language models such as CLIP exhibit remarkable zero-shot generalization, yet their vulnerability to adversarial perturbations constitutes a critical obstacle to deployment in security-sensitive settings. Under standard projected gradient descent (PGD) attacks with $L_\infty$ budget $\epsilon=4/255$, zero-shot classification accuracy collapses from above $80\%$ to below $5\%$. Test-Time Counterattack (TTC, CVPR 2025) introduced a training-free defense paradigm in which a corrective perturbation is optimized at inference time via PGD to realign adversarial embeddings with their clean counterparts. However, the standard gradient-based optimizer employed by TTC converges slowly under the tight iteration budgets ($K=2$--$4$) imposed by inference latency constraints. In this paper, we propose Nesterov Accelerated Counterattack (NAC), which replaces the standard PGD update with Nesterov accelerated gradient (NAG). The modification consists of a single principled change: gradients are evaluated at a look-ahead position $\delta_t + \mu v_t$ rather than at the current iterate $\delta_t$. While NAG's optimal $O(1/k^2)$ convergence rate is provably guaranteed only in the convex setting, we demonstrate that its empirical acceleration transfers effectively to the non-convex counterattack landscape, yielding substantial practical gains at identical computational cost. Extensive experiments across six benchmark datasets demonstrate that NAC consistently outperforms TTC under PGD attacks at three perturbation strengths ($\epsilon=1/255, 2/255, 4/255$) and under the AutoAttack ensemble. On CIFAR-10 at $\epsilon=4/255$, NAC with two counterattack steps improves robust accuracy from $6.98\%$ to $16.61\%$ ($+9.63$ percentage points); with four steps, NAC reaches $37.05\%$ compared to TTC's $25.37\%$. Controlled ablation experiments isolate Nesterov's look-ahead mechanism from standard momentum accumulation, confirming that the look-ahead---not momentum alone---accounts for the observed gains. NAC further generalizes across ViT architectures and yields complementary improvements when combined with adversarially fine-tuned models.
}

\keywords{Nesterov accelerated gradient, adversarial robustness, vision-language model, CLIP, test-time defense, counterattack}

\maketitle

% ============================================================
\section{Introduction}\label{sec:intro}

Vision-language models pretrained on web-scale image-text pairs, most notably CLIP~\cite{radford2021clip}, have demonstrated remarkable zero-shot generalization across diverse visual recognition tasks. By learning aligned representations from natural language supervision, these models circumvent the need for task-specific training data and have become foundational components in a broad range of downstream applications. However, a substantial body of work has established that CLIP inherits the well-known vulnerability of deep neural networks to adversarial perturbations~\cite{goodfellow2015explaining,madry2018towards,croce2020reliable}. Concretely, an adversary can reduce CLIP's zero-shot classification accuracy from above $80\%$ to below $5\%$ by adding imperceptible $L_\infty$-bounded noise with a budget as small as $\epsilon = 4/255$~\cite{mao2023understanding}. This brittleness raises significant concerns for the deployment of vision-language models in security-sensitive settings.

Existing defenses against adversarial attacks on vision-language models can be broadly categorized into two paradigms. The first, adversarial fine-tuning (AFT), modifies model weights through adversarial training~\cite{mao2023understanding,schlarmann2024robust,xu2024pretrained,chen2025advclip}. While AFT methods have achieved non-trivial improvements in robust accuracy, they incur a well-documented trade-off: the adversarial training objective systematically degrades the zero-shot generalization capability that constitutes CLIP's primary advantage. For instance, TeCoA~\cite{mao2023understanding} reported that standard adversarial training reduced clean accuracy from $64.56\%$ to $18.49\%$ before its proposed text-guided loss was introduced to partially recover the lost generalization. The second paradigm, test-time defense, operates entirely at inference without modifying model parameters~\cite{xing2025clip,jiang2026diversifying,liu2026directional,nie2025contrastive,li2026agc}. By design, these methods preserve the model's zero-shot capabilities and can be applied in a plug-and-play manner.

Among test-time approaches, Test-Time Counterattack (TTC)~\cite{xing2025clip} (CVPR 2025) introduced a defense paradigm that departs fundamentally from prior strategies: rather than passively detecting or purifying adversarial perturbations, TTC actively \emph{counterattacks}---it employs projected gradient descent (PGD) to optimize a corrective perturbation $\delta$ that pushes the adversarial image's embedding back toward the clean image's embedding in CLIP's representation space. However, TTC inherits a critical limitation from its PGD backbone: the $O(1/k)$ sublinear convergence rate of gradient descent~\cite{nesterov1983method}. Under inference-time constraints, counterattack is allotted merely $K=2$--$4$ iterations, at which point the optimization remains severely under-converged. The empirical evidence is revealing: increasing TTC's counterattack steps from $K=2$ to $K=4$ lifts robust accuracy on CIFAR-10 from $7.07\%$ to $25.37\%$ under PGD $\epsilon = 4/255$---an $18.30$ percentage point gap that confirms substantial untapped potential remains at small iteration budgets.

We identify the choice of optimizer as a critical yet previously unexamined design dimension in test-time counterattack. Motivated by this observation, we propose \textbf{N}esterov \textbf{A}ccelerated \textbf{C}ounterattack (NAC), which replaces the standard PGD update with Nesterov accelerated gradient (NAG)~\cite{nesterov1983method}. The Nesterov formulation introduces a single, principled modification: gradients are evaluated at a \textbf{look-ahead} position $\delta_t + \mu v_t$---the extrapolated future state obtained by applying the current velocity---rather than at the current position $\delta_t$. Under locally smooth conditions, this adjustment yields the optimal $O(1/k^2)$ convergence rate for first-order methods in the convex setting~\cite{nesterov1983method}, a quadratic improvement over PGD's $O(1/k)$. The modification is minimal in implementation---requiring a change only in where the gradient is evaluated---and incurs \emph{no additional computational cost}: the forward and backward passes are identical in number and complexity to those of TTC.

The practical effectiveness of NAC is enabled by an empirical property of CLIP's ViT embedding space: under the small perturbation budget $\epsilon = 4/255 \approx 0.016$, the optimization landscape is locally well-behaved, allowing the look-ahead gradient to serve as a substantially more reliable descent direction. Figure~\ref{fig:framework} provides a conceptual comparison of the two gradient computation schemes. Quantitatively, NAC with two counterattack steps (NAC-2) achieves $16.61\%$ robust accuracy on CIFAR-10 under PGD $\epsilon=4/255$, compared to TTC-2's $6.98\%$---a gain of $9.63$ percentage points. With four steps, NAC-4 reaches $37.05\%$ versus TTC-4's $25.37\%$ ($+11.68$ percentage points).

Our contributions are as follows.
\begin{enumerate}[leftmargin=*]
\item We identify optimizer selection as a critical yet overlooked factor in test-time counterattack, and propose NAC---a Nesterov modification of TTC that yields substantially faster convergence at identical computational cost.
\item We analyze the convergence behavior of NAG in the context of test-time counterattack, elucidating why the look-ahead mechanism is particularly effective in CLIP's locally smooth embedding space under small perturbation budgets.
\item Through extensive experiments spanning six datasets, three attack strengths, and the AutoAttack ensemble, we demonstrate that NAC consistently outperforms TTC across all evaluated settings.
\item Through controlled ablation studies, we isolate the contribution of Nesterov's look-ahead mechanism from that of standard momentum accumulation, establishing that the look-ahead---not momentum alone---accounts for the observed gains.
\end{enumerate}

\begin{figure}[ht]\centering
\includegraphics[width=0.95\textwidth]{fig1_framework.pdf}
\caption{Conceptual comparison of TTC and NAC gradient computation. TTC evaluates the gradient at the current perturbation $\delta_t$, yielding a locally steepest descent direction. NAC evaluates the gradient at the look-ahead position $\delta_t + \mu v_t$, which, under local smoothness, provides a superior approximation of the gradient at the next iterate, enabling accelerated convergence.}\label{fig:framework}
\end{figure}

% ============================================================
\section{Related Work}\label{sec:related}

\subsection{Adversarial Robustness of Vision-Language Models}

The adversarial vulnerability of CLIP was first systematically characterized by Mao et al.~\cite{mao2023understanding}, who established that zero-shot CLIP suffers severe accuracy degradation under both $L_\infty$ and $L_2$ bounded perturbations. This finding motivated a growing body of work on defenses specifically tailored to vision-language models.

\emph{Adversarial fine-tuning (AFT).} TeCoA~\cite{mao2023understanding} (CVPR 2023) introduced text-guided contrastive adversarial training, achieving $38.18\%$ robust accuracy at the cost of an $8.59$ percentage point drop in clean accuracy relative to the standard model. FARE~\cite{schlarmann2024robust} (ICML 2024 Oral) proposed unsupervised adversarial fine-tuning of the vision encoder, avoiding the need for labeled data and achieving $67.0\%$ clean accuracy with $19.6\%$ robust accuracy. PMG-AFT~\cite{xu2024pretrained} (CVPR 2024) leveraged a frozen pretrained model as a guidance teacher, attaining $46.60\%$ clean and $32.51\%$ robust accuracy. AdvCLIP-LoRA~\cite{chen2025advclip} explored parameter-efficient fine-tuning via low-rank adaptation, substantially reducing the number of trainable parameters. A common limitation across all AFT methods is the inherent tension between adversarial robustness and zero-shot generalization~\cite{gowal2021improving}: weight modification during training inevitably shifts the model away from the pretrained representation that enables zero-shot transfer.

\emph{Attention-based and prompt-based defenses.} Orthogonal to weight-space methods, TGA-ZSR~\cite{yu2024text} (NeurIPS 2024) proposed text-guided attention constraints that suppress spurious attention patterns induced by adversarial perturbations, achieving $42.09\%$ robust and $56.44\%$ clean accuracy. Prompt-tuning approaches such as TAPT~\cite{wang2025tapt} learn defensive soft prompts at test time, avoiding permanent modification of the text encoder. NAC is complementary to both attention-based and prompt-based defenses, as it operates solely on the counterattack optimization loop and can be combined with these orthogonal strategies.

\subsection{Test-Time Defenses for CLIP}

Test-time defenses have emerged as an attractive paradigm due to their training-free nature and compatibility with any pretrained model. TTC~\cite{xing2025clip} (CVPR 2025) pioneered the counterattack framework, optimizing a corrective perturbation via PGD to minimize the $L_2$ distance between adversarial and clean embeddings in CLIP's feature space. Across 16 datasets, TTC raised average robust accuracy from $3.11\%$ to $20.64\%$. DOC~\cite{jiang2026diversifying} (AAAI 2026 Oral) extended TTC through four complementary mechanisms: orthogonal gradient directions for exploration diversity, directional sensitivity analysis, sigmoid-based learnable tau gating, and gradient normalization. DBD~\cite{liu2026directional} (ICLR 2026) exploited the directional bias inherent in adversarial feature shifts, demonstrating that the attack direction itself contains information exploitable for defense. CSR~\cite{nie2025contrastive} proposed frequency-domain spectral detection combined with PGD-based purification, operating without training. AGC~\cite{li2026agc} employed geodesic correction on the embedding manifold. TTP~\cite{li2026ttp} explored test-time padding for simultaneous adversarial detection and robust adaptation.

Despite the diversity of these approaches, a shared characteristic unites them: every gradient-based test-time defense, including TTC, DOC, DBD, CSR, and AGC, employs standard PGD or simple Polyak momentum as its optimizer. The choice of optimizer has received little systematic examination. Our work identifies and exploits this overlooked design dimension, demonstrating that a principled optimizer substitution---replacing PGD with Nesterov accelerated gradient---yields consistent and substantial gains at no additional computational cost.

\subsection{Nesterov Accelerated Gradient in Deep Learning}

Nesterov's accelerated gradient method~\cite{nesterov1983method} achieves the optimal $O(1/k^2)$ convergence rate for convex smooth optimization, a quadratic improvement over gradient descent's $O(1/k)$. Its importance in deep learning was established by Sutskever et al.~\cite{sutskever2013importance}, who demonstrated that Nesterov momentum consistently outperforms standard Polyak momentum for training deep neural networks. NAG's principles were subsequently incorporated into adaptive methods such as Adam~\cite{kingma2015adam} (via NAdam). In the context of adversarial robustness, NAG has been applied exclusively to \textbf{attack generation}: NI-FGSM~\cite{dong2018boosting} (CVPR 2018) employed Nesterov momentum to enhance the transferability of adversarial examples across models, and Lin et al.~\cite{lin2024nesterov} (NeurIPS 2024) further combined NAG with scale invariance for stronger transfer-based attacks. To our knowledge, Nesterov acceleration has not previously been applied to adversarial defense or counterattack. Our work bridges this gap, repurposing a well-established optimization technique from attack generation to test-time defense.

% ============================================================
\section{Method}\label{sec:method}

\subsection{Preliminaries}

Let $f_I: \mathbb{R}^{3 \times H \times W} \to \mathbb{R}^d$ denote CLIP's vision encoder and $f_T: \mathcal{V}^L \to \mathbb{R}^d$ its text encoder. Given a clean image $x$ and a set of class names $\{c_1, \dots, c_C\}$, zero-shot classification proceeds by computing cosine similarities between the image embedding $f_I(x)$ and text embeddings $f_T(t_i)$ obtained from prompt templates $t_i$ (e.g., ``a photo of a $\{c_j\}$''), then applying softmax over the similarity scores.

An adversarial example $x_{\text{adv}} = x + \delta_{\text{att}}$ is generated by maximizing the classification loss subject to an $L_\infty$ perturbation budget $\epsilon_{\text{att}}$:
\begin{equation}
\delta_{\text{att}} = \arg\max_{\|\delta\|_\infty \leq \epsilon_{\text{att}}} \mathcal{L}_{\text{cls}}\big(f_I(x+\delta), y\big) \label{eq:attack}
\end{equation}
where $\mathcal{L}_{\text{cls}}$ denotes the cross-entropy loss and $y$ the ground-truth label. In practice, Equation~\ref{eq:attack} is solved via projected gradient descent (PGD)~\cite{madry2018towards}.

\subsection{Test-Time Counterattack (TTC)}

TTC~\cite{xing2025clip} reformulates defense as an optimization problem in the embedding space. Rather than modifying the classifier boundary, TTC seeks a corrective perturbation $\delta$ that aligns the adversarial embedding with the clean embedding:
\begin{equation}
\min_{\|\delta\|_\infty \leq \epsilon_{\text{ttc}}} \mathcal{L}(\delta) = \big\|f_I(x_{\text{adv}} + \delta) - f_I(x)\big\|^2_2 \label{eq:ttc_obj}
\end{equation}
where $\epsilon_{\text{ttc}}$ denotes the counterattack perturbation budget. This objective is optimized via $K$ steps of signed PGD:
\begin{equation}
\delta_{t+1} = \Pi_{\epsilon_{\text{ttc}}}\!\big(\delta_t - \alpha \cdot \operatorname{sign}(\nabla \mathcal{L}(\delta_t))\big) \label{eq:pgd_update}
\end{equation}
where $\Pi_{\epsilon}$ denotes projection onto the $L_\infty$ ball of radius $\epsilon$, and $\alpha$ is the step size. After $K$ steps, a tau-threshold gating mechanism produces a weighted combination of the intermediate iterates $\{\delta_0, \delta_1, \dots, \delta_K\}$ to form the final counter-perturbation $\Delta$.

\subsection{NAC: Nesterov Accelerated Counterattack}

The core limitation of TTC lies in the convergence rate of the underlying PGD optimizer. Standard gradient descent achieves $O(1/K)$ convergence for smooth objectives~\cite{nesterov1983method}, meaning that under tight iteration budgets ($K=2$--$4$), the optimization remains far from its stationary point. NAC addresses this by replacing the PGD update with Nesterov accelerated gradient (NAG). The NAC update rules are:
\begin{align}
v_{t+1} &= \mu \cdot v_t + \alpha \cdot \operatorname{sign}\!\big(\nabla \mathcal{L}(\delta_t + \mu \cdot v_t)\big) \label{eq:nac_vel} \\
\delta_{t+1} &= \Pi_{\epsilon_{\text{ttc}}}\!\big(\delta_t + v_{t+1}\big) \label{eq:nac_update}
\end{align}
where $v_t$ is the velocity (accumulated momentum) and $\mu \in [0, 1)$ is the momentum coefficient. The distinguishing characteristic of NAC, embodied in Equation~\ref{eq:nac_vel}, is the \textbf{look-ahead gradient evaluation}: the gradient is computed at the extrapolated position $\delta_t + \mu \cdot v_t$ rather than at the current iterate $\delta_t$. Algorithm~\ref{alg:nac} presents the complete pseudocode.

\begin{algorithm}[ht]
\caption{NAC: Nesterov Accelerated Counterattack}\label{alg:nac}
\begin{algorithmic}[1]
\Require CLIP vision encoder $f_I$, adversarial image $x_{\text{adv}}$, clean embedding $e_{\text{clean}}=f_I(x)$, steps $K$, step size $\alpha$, momentum $\mu$, budget $\epsilon$, tau threshold $\tau$, temperature $\beta$
\Ensure Counter-perturbation $\Delta$
\State $\delta \leftarrow 0$, $v \leftarrow 0$, $D \leftarrow [\delta]$
\For{$t = 1$ to $K$}
    \State $x_{\text{la}} \leftarrow x_{\text{adv}} + \delta + \mu \cdot v$ \Comment{Look-ahead position (NAC)}
    \State $\text{loss} \leftarrow \big\|f_I(x_{\text{la}}) - e_{\text{clean}}\big\|^2_2$
    \State $g \leftarrow \nabla_\delta \,\text{loss}$
    \State $v \leftarrow \mu \cdot v + \alpha \cdot \operatorname{sign}(g)$ \Comment{Nesterov velocity update}
    \State $\delta \leftarrow \Pi_{\epsilon}(\delta + v)$
    \State $D.\text{append}(\delta)$
\EndFor
\State $\Delta \leftarrow \textsc{TauWeightedAverage}(D, \tau, \beta)$
\State \Return $\Delta$
\end{algorithmic}
\end{algorithm}

\subsection{Convergence Analysis}

To understand the advantage of NAC over TTC, we analyze the convergence behavior of the underlying optimizers. Let $g(\delta) = \|f_I(x_{\text{adv}} + \delta) - f_I(x)\|^2_2$ denote the counterattack objective.

\begin{proposition}[Convergence of gradient descent]\label{prop:pgd}
Assume $g$ is $L$-smooth. Then PGD with step size $\alpha = 1/L$ satisfies:
\begin{equation}
\min_{t < K} \|\nabla g(\delta_t)\|^2 \leq \frac{2L\,(g(\delta_0) - g^*)}{K} = O(1/K).
\end{equation}
\end{proposition}

\begin{proposition}[Convergence of Nesterov accelerated gradient]\label{prop:nag}
Assume $g$ is $L$-smooth and convex. Then NAG with step size $\alpha = 1/L$ satisfies:
\begin{equation}
g(\delta_K) - g^* \leq \frac{2L\,\|\delta_0 - \delta^*\|^2}{K^2} = O(1/K^2).
\end{equation}
\end{proposition}

\textbf{Discussion.} We acknowledge that the counterattack objective $g(\delta) = \|f_I(x_{\text{adv}} + \delta) - f_I(x)\|^2$ is, in general, non-convex in $\delta$ due to the nonlinearity of the vision encoder $f_I$, and a rigorous convergence analysis of NAG in the non-convex setting remains an active area of research. Ghadimi and Lan~\cite{ghadimi2016accelerated} established that randomized NAG achieves an $O(1/K^2)$ rate for finding an $\epsilon$-approximate stationary point in non-convex stochastic programming, but their analysis requires averaging over multiple runs and additional variance reduction techniques that are impractical under a $K=2$--$4$ iteration budget. More recent work on non-convex accelerated methods has focused on adaptive restart schemes and negative momentum variants, which share the high-level principle of look-ahead exploration with NAC but are designed for large-scale training rather than single-instance test-time optimization.

Nevertheless, several considerations motivate the practical effectiveness of NAC despite the absence of a formal non-convex guarantee. First, the counterattack budget $\epsilon_{\text{ttc}} = 4/255 \approx 0.016$ constrains optimization to a highly local neighborhood of the adversarial image. Within such a restricted region, the embedding function $f_I$ is empirically observed to be approximately smooth, as evidenced by the monotonic improvement of both TTC and NAC with additional counterattack steps. Second, even without global convexity, NAG's look-ahead mechanism provides a strictly better gradient approximation than the current-position gradient whenever the objective admits a local quadratic expansion---a condition satisfied by any neural network with twice-differentiable activations in a sufficiently small neighborhood. Under a second-order Taylor expansion $g(\delta + \Delta) \approx g(\delta) + \nabla g(\delta)^\top \Delta + \frac{1}{2} \Delta^\top H(\delta) \Delta$, the NAG update with look-ahead $\delta + \mu v$ can be interpreted as preconditioning the gradient by an implicit approximation to the inverse Hessian, effectively increasing the step size along low-curvature directions~\cite{nesterov1983method}. Third, the practical distinction between $O(1/K)$ and $O(1/K^2)$ is most pronounced in the small-$K$ regime ($K=2$--$4$) that characterizes test-time defense, where asymptotic constants dominate and each iteration's gradient quality directly impacts final accuracy. We leave a formal non-convex convergence analysis of NAC---potentially along the lines of Ghadimi and Lan's randomized scheme adapted to the single-instance, small-$K$ setting---as an open problem for future theoretical work.

\subsection{Distinction from Standard Momentum}

We contrast NAC's Nesterov update with standard (Polyak) momentum to clarify the distinction. The Polyak heavy-ball update is:
\begin{align}
m_{t+1} &= \mu \cdot m_t + \alpha \cdot \operatorname{sign}\!\big(\nabla g(\delta_t)\big) \label{eq:polyak} \\
\delta_{t+1} &= \Pi_{\epsilon}\!\big(\delta_t + m_{t+1}\big)
\end{align}
where the gradient is evaluated at the \textbf{current} position $\delta_t$. In contrast, NAC evaluates the gradient at the \textbf{look-ahead} position $\delta_t + \mu \cdot v_t$. The distinction is fundamental rather than cosmetic: standard momentum smooths past gradients to dampen oscillations, whereas Nesterov momentum predicts the future gradient direction by evaluating at the extrapolated state. As we demonstrate empirically in Section~\ref{sec:ablation}, this distinction has substantial practical consequences: the look-ahead mechanism contributes the majority of NAC's improvement over TTC, far exceeding the benefit of momentum accumulation alone.

% ============================================================
\section{Experiments}\label{sec:experiments}

\subsection{Experimental Setup}

\textbf{Datasets.} We evaluate on six benchmark datasets spanning diverse visual domains, as summarized in Table~\ref{tab:datasets}: CIFAR-10 and CIFAR-100~\cite{krizhevsky2009learning} (general object recognition, $32\times32$ resolution), STL-10~\cite{coates2011analysis} (general objects, $96\times96$), Flowers-102~\cite{nilsback2008automated} (fine-grained flower classification), Describable Textures Dataset (DTD)~\cite{cimpoi2014describing} (texture recognition), and ImageNet-100 (a 100-class subset of ImageNet~\cite{deng2009imagenet}, ${\sim}500$ resolution). All datasets are evaluated in the zero-shot setting; no training or fine-tuning is performed on any dataset.

\textbf{Models.} Unless otherwise specified, all experiments use CLIP ViT-B/32~\cite{radford2021clip} as the default model. Cross-architecture experiments additionally evaluate ViT-B/16 (finer patch size, more visual tokens) and RN50 (ResNet-50 backbone).

\textbf{Attack configuration.} Following the standard evaluation protocol of TTC~\cite{xing2025clip}, we employ PGD under $L_\infty$ norm with $\epsilon_{\text{att}} \in \{1/255, 2/255, 4/255\}$, $10$ attack steps, and step size $1/255$. For stronger evaluation, we additionally employ AutoAttack~\cite{croce2020reliable} with the standard APGD-CE + APGD-DLR ensemble.

\textbf{Counterattack configuration.} For both TTC and NAC, we adopt the default TTC hyperparameters: counterattack budget $\epsilon_{\text{ttc}} = 4/255$, step size $\alpha = 1/255$, tau threshold $\tau = 0.2$, and temperature $\beta = 2.0$. NAC additionally uses momentum coefficient $\mu = 0.9$. The default number of counterattack steps is $K=2$ for main experiments; $K=4$ is evaluated in the N-step scaling analysis. All results are averaged over three random seeds ($0, 1, 2$); the maximum observed standard deviation is below $0.2$ percentage points, confirming stable convergence.

\begin{table}[ht]\centering\small
\caption{Dataset statistics and characteristics. The six datasets span general object recognition (CIFAR-10, CIFAR-100, STL-10, ImageNet-100), fine-grained classification (Flowers-102), and texture analysis (DTD), providing a comprehensive evaluation of robustness across diverse visual domains.}\label{tab:datasets}
\begin{tabular}{@{}lrrll@{}}
\toprule Dataset & Classes & Images & Resolution & Domain \\
\midrule CIFAR-10 & 10 & 10,000 & $32\times32$ & General objects \\
CIFAR-100 & 100 & 10,000 & $32\times32$ & General objects \\
STL-10 & 10 & 8,000 & $96\times96$ & General objects \\
Flowers-102 & 102 & 6,149 & ${\sim}500$ & Fine-grained \\
DTD & 47 & 1,880 & ${\sim}300$ & Textures \\
ImageNet-100 & 100 & 5,000 & ${\sim}500$ & General objects \\
\bottomrule
\end{tabular}
\end{table}

\subsection{Main Results}

Table~\ref{tab:main} reports robust accuracy under PGD $\epsilon_{\text{att}}=1/255$ with the default 2-step counterattack. NAC consistently outperforms TTC across all six datasets, with gains ranging from $+2.19$ percentage points on CIFAR-100 to $+7.14$ percentage points on Flowers-102. The average improvement of $+3.83$ percentage points is substantial in the context of test-time defense, where gains are typically modest due to the fixed model and limited inference budget.

Notably, the magnitude of NAC's improvement varies considerably across datasets. The largest gain is observed on Flowers-102 ($+7.14$pp), a fine-grained dataset where subtle inter-class distinctions are easily disrupted by adversarial perturbations---and where the accelerated convergence of NAG appears particularly beneficial for recovering discriminative embedding structure. In contrast, the more modest gain on CIFAR-100 ($+2.19$pp) likely reflects the higher difficulty of the 100-class setting under zero-shot evaluation, where the absolute robust accuracy ($16.87\%$) leaves less headroom for any defense method.

\begin{table}[ht]\centering\small
\caption{Robust accuracy (\%) under PGD attack with $\epsilon_{\text{att}}=1/255$, using 2-step counterattack ($K=2$). NAC consistently outperforms TTC across all six datasets, with an average gain of $+3.83$ percentage points. Bold indicates best result.}\label{tab:main}
\begin{tabular}{@{}lrrrr@{}}
\toprule Dataset & Clean & TTC & NAC & Gain (pp) \\
\midrule CIFAR-10 & 84.89 & 27.95 & \textbf{30.94} & +2.99 \\
CIFAR-100 & 58.19 & 14.68 & \textbf{16.87} & +2.19 \\
STL-10 & 96.67 & 77.11 & \textbf{80.33} & +3.22 \\
Flowers-102 & 62.73 & 38.62 & \textbf{45.76} & +7.14 \\
DTD & 42.71 & 26.54 & \textbf{29.68} & +3.14 \\
ImageNet-100 & 67.82 & 3.68 & \textbf{7.99} & +4.31 \\
\midrule Average & -- & 31.43 & \textbf{35.26} & +3.83 \\
\bottomrule
\end{tabular}
\end{table}

\subsection{Multi-Attack Strength Evaluation}

To assess the robustness of NAC's advantage across varying threat intensities, we evaluate under three attack budgets $\epsilon_{\text{att}} \in \{1/255, 2/255, 4/255\}$. Table~\ref{tab:multi} reports average robust accuracy across all six datasets. NAC delivers consistent gains at every attack strength, with the largest absolute improvement at $\epsilon_{\text{att}}=2/255$ ($+7.49$pp). Notably, the gain at $\epsilon_{\text{att}}=4/255$ ($+5.49$pp) is smaller than at $\epsilon_{\text{att}}=2/255$, a counterintuitive pattern that warrants explanation. We attribute this to the severity of the adversarial perturbation at $\epsilon_{\text{att}}=4/255$: when the attack budget is large, the adversarial image's embedding is displaced so far from the clean embedding that the counterattack objective $g(\delta) = \|f_I(x_{\text{adv}} + \delta) - f_I(x)\|^2$ loses its local smoothness. The embedding function $f_I$ becomes highly nonlinear over the required traversal distance, and the look-ahead gradient---which relies on local regularity---provides a less reliable descent direction. Under the strongest attack ($\epsilon_{\text{att}}=4/255$), where PGD reduces clean accuracy to below $5\%$, NAC nevertheless more than doubles the average robust accuracy achieved by TTC ($10.89\%$ vs.\ $5.40\%$), underscoring the practical value of accelerated convergence even in partially degraded optimization landscapes.

\begin{table}[ht]\centering\small
\caption{Average robust accuracy (\%) across six datasets under PGD with varying attack budgets $\epsilon_{\text{att}}$. NAC's advantage is consistent across all threat levels, with the gap widening at intermediate perturbation strength ($\epsilon=2/255$).}\label{tab:multi}
\begin{tabular}{@{}lrrr@{}}
\toprule $\epsilon_{\text{attack}}$ & TTC & NAC & Gain (pp) \\
\midrule 1/255 & 31.43 & \textbf{35.26} & +3.83 \\
2/255 & 22.07 & \textbf{29.56} & +7.49 \\
4/255 & 5.40 & \textbf{10.89} & +5.49 \\
\bottomrule
\end{tabular}
\end{table}

\subsection{N-Step Scaling Analysis}

A key prediction of our convergence analysis is that NAC should extract more value per additional iteration than TTC. Table~\ref{tab:nstep} tests this prediction by evaluating both methods with $K=2$ and $K=4$ counterattack steps under the strongest attack ($\epsilon_{\text{att}}=4/255$) on CIFAR-10 and STL-10.\footnotemark[1] The results are consistent with our analysis: NAC-2 achieves robust accuracy that substantially exceeds TTC-2 on both datasets, and the gap widens further at $K=4$. The per-step efficiency gain is notable: NAC with 2 steps achieves $16.47\%$ on CIFAR-10 and $34.30\%$ on STL-10, while TTC requires 4 steps to reach $25.37\%$ and $48.29\%$, respectively---NAC-2 already recovers $65\%$ (CIFAR-10) and $71\%$ (STL-10) of TTC-4's accuracy at half the computational cost. This efficiency directly benefits latency-constrained deployment scenarios where each additional iteration incurs non-trivial overhead.
\footnotetext[1]{Minor numerical variations across tables (e.g., NAC-2: $16.47\%$ in Table~\ref{tab:nstep} vs.\ $16.61\%$ in Table~\ref{tab:ablation}; TTC-2: $6.98\%$ in Table~\ref{tab:ablation} vs.\ $7.07\%$ in Table~\ref{tab:nstep}) arise from different random seeds between experimental runs. All are within the $\pm 0.2$pp stability range reported in Section~\ref{sec:experiments}.}

\begin{table}[ht]\centering\small
\caption{Robust accuracy (\%) with varying counterattack steps under PGD $\epsilon_{\text{att}}=4/255$. NAC consistently extracts more value per additional iteration than TTC, consistent with the $O(1/K^2)$ vs.\ $O(1/K)$ convergence distinction.}\label{tab:nstep}
\begin{tabular}{@{}lrlr@{}}
\toprule Method & Steps & CIFAR-10 & STL-10 \\
\midrule TTC & 2 & 7.07 & 17.44 \\
NAC & 2 & \textbf{16.47} & \textbf{34.30} \\
TTC & 4 & 25.37 & 48.29 \\
NAC & 4 & \textbf{37.05} & \textbf{66.74} \\
\bottomrule
\end{tabular}
\end{table}

\subsection{Ablation: Isolating the Look-Ahead Mechanism}\label{sec:ablation}

A natural question is whether NAC's improvement originates from Nesterov's look-ahead gradient computation or merely from the introduction of momentum---a standard technique known to accelerate optimization. To answer this question, we conduct a controlled ablation comparing three configurations under identical hyperparameter settings: (1) TTC with no momentum, (2) TTC augmented with standard Polyak momentum ($\mu=0.9$, gradient at current position), and (3) NAC (Nesterov look-ahead, $\mu=0.9$).

Table~\ref{tab:ablation} presents the results. Standard momentum provides a modest improvement over TTC ($+1.48$pp on CIFAR-10, $+3.88$pp on STL-10), confirming that momentum alone offers some benefit. However, NAC's Nesterov formulation yields an additional $+8.15$pp on CIFAR-10 and $+12.99$pp on STL-10 beyond standard momentum---gains approximately $5.5\times$ and $3.3\times$ larger, respectively. This pronounced asymmetry isolates the look-ahead gradient evaluation as the primary driver of NAC's effectiveness. Figure~\ref{fig:ablation} provides a visual comparison of the three configurations.

\begin{table}[ht]\centering\small
\caption{Ablation study isolating the effect of Nesterov's look-ahead mechanism from standard momentum accumulation. All configurations use identical hyperparameters ($\mu=0.9$, $\epsilon_{\text{att}}=4/255$, $K=2$). Standard momentum contributes a modest gain; the look-ahead mechanism contributes the majority of NAC's improvement.}\label{tab:ablation}
\begin{tabular}{@{}lrr@{}}
\toprule Method & CIFAR-10 & STL-10 \\
\midrule TTC (no momentum) & 6.98 & 17.32 \\
$+$ Standard momentum & 8.46 & 21.20 \\
$+$ Nesterov look-ahead (NAC) & \textbf{16.61} & \textbf{34.19} \\
\bottomrule
\end{tabular}
\end{table}

\begin{figure}[ht]\centering
\includegraphics[width=0.85\textwidth]{fig2_ablation.pdf}
\caption{Ablation comparison of TTC (no momentum), TTC with standard Polyak momentum, and NAC (Nesterov look-ahead). The Nesterov look-ahead contributes substantially more improvement than standard momentum alone, confirming the distinct and complementary nature of the two mechanisms.}\label{fig:ablation}
\end{figure}

\subsection{Momentum Coefficient Sensitivity}

Table~\ref{tab:mu} reports NAC performance as a function of the momentum coefficient $\mu$, with $\mu=0$ corresponding to standard TTC. Robust accuracy improves monotonically with $\mu$, from $6.06\%$ at $\mu=0$ to $20.54\%$ at $\mu=0.9$, representing a $14.48$ percentage point gain. Further increasing $\mu$ to $0.99$ yields a marginal additional improvement ($21.46\%$), suggesting diminishing returns beyond $\mu=0.9$. We adopt $\mu=0.9$ as the default value, consistent with prior work on NAG in deep learning~\cite{nesterov1983method,sutskever2013importance}. The monotonic relationship between $\mu$ and robust accuracy (visualized in Figure~\ref{fig:mu}) provides further evidence that NAC's improvement is a genuine consequence of accelerated optimization rather than an artifact of hyperparameter tuning.

\begin{table}[ht]\centering\small
\caption{NAC robust accuracy (\%) as a function of momentum coefficient $\mu$ on CIFAR-10 under PGD $\epsilon_{\text{att}}=4/255$ with $K=5$ steps. Accuracy improves monotonically with $\mu$, confirming the systematic benefit of look-ahead gradient evaluation.}\label{tab:mu}
\begin{tabular}{@{}lcccccc@{}}
\toprule $\mu$ & 0 (TTC) & 0.1 & 0.5 & 0.7 & 0.9 & 0.99 \\
\midrule NAC & 6.06 & 10.73 & 15.81 & 18.16 & \textbf{20.54} & 21.46 \\
\bottomrule
\end{tabular}
\end{table}

\begin{figure}[ht]\centering
\includegraphics[width=0.7\textwidth]{fig3_mu.pdf}
\caption{Robust accuracy of NAC as a function of the momentum coefficient $\mu$. The monotonic trend confirms that the look-ahead mechanism provides systematic benefit, with the gain saturating near $\mu=0.9$.}\label{fig:mu}
\end{figure}

\subsection{Evaluation under AutoAttack}

To assess NAC's robustness against a stronger, ensemble-based attack, we employ AutoAttack~\cite{croce2020reliable} with the standard APGD-CE + APGD-DLR combination. Table~\ref{tab:aa} reports results on CIFAR-10 at $\epsilon_{\text{att}}=4/255$. NAC maintains a clear advantage over TTC ($+3.73$pp), confirming that the benefit of Nesterov acceleration is not specific to PGD-based attacks and generalizes to stronger adaptive adversaries. This result is particularly important given the well-known phenomenon of gradient masking, where defenses that appear effective against PGD can fail against more thorough attacks. The consistent improvement under AutoAttack suggests that NAC represents a genuine robustness gain rather than a gradient obfuscation effect.

\begin{table}[ht]\centering\small
\caption{Robust accuracy (\%) under AutoAttack (APGD-CE + APGD-DLR) on CIFAR-10, $\epsilon_{\text{att}}=4/255$. NAC's advantage persists under the stronger ensemble attack, ruling out gradient obfuscation as an explanation for the observed gains.}\label{tab:aa}
\begin{tabular}{@{}lrrrr@{}}
\toprule Method & Clean & Adv & Defended & Gain (pp) \\
\midrule TTC & 84.90 & 0.05 & 7.56 & -- \\
NAC & 84.90 & 0.05 & \textbf{11.29} & +3.73 \\
\bottomrule
\end{tabular}
\end{table}

\subsection{Cross-Architecture Generalization}

We evaluate NAC on ViT-B/16, which employs finer patch partitioning ($16\times16$ vs.\ $32\times32$) and consequently processes more visual tokens, resulting in a different embedding geometry. Table~\ref{tab:cross} reports results on CIFAR-10 at $\epsilon_{\text{att}}=4/255$. NAC generalizes to ViT-B/16, delivering a $+6.15$pp improvement over TTC. The slightly smaller gain compared to ViT-B/32 ($+9.63$pp) likely reflects the higher baseline robust accuracy of ViT-B/16 ($8.51\%$ vs.\ $6.98\%$), which reduces the headroom for improvement. On RN50, both TTC and NAC recover negligible robust accuracy; we attribute this to the fundamentally different embedding geometry of ResNet's attention-pooled features, for which the counterattack hyperparameters (tuned on ViT) are suboptimal. We note that this limitation is shared with TTC~\cite{xing2025clip} and leave RN50-specific hyperparameter optimization to future work.

\begin{table}[ht]\centering\small
\caption{Cross-architecture evaluation on CIFAR-10 under PGD $\epsilon_{\text{att}}=4/255$. NAC generalizes to ViT-B/16, confirming that the Nesterov acceleration benefit is not tied to a specific patch size or token count. RN50 results reflect ViT-tuned hyperparameters.}\label{tab:cross}
\begin{tabular}{@{}lrrrr@{}}
\toprule Architecture & Clean & TTC & NAC & Gain (pp) \\
\midrule ViT-B/32 & 84.89 & 6.98 & \textbf{16.61} & +9.63 \\
ViT-B/16 & 87.24 & 8.51 & \textbf{14.66} & +6.15 \\
\bottomrule
\end{tabular}
\end{table}

\subsection{Complementary Gains with Adversarial Fine-Tuning}

Test-time counterattacks are orthogonal to adversarial fine-tuning: AFT modifies model weights during training, while NAC operates post hoc at inference without weight modification. To assess complementarity, we evaluate NAC on top of two representative AFT models, TeCoA~\cite{mao2023understanding} and FARE~\cite{schlarmann2024robust}, under the strongest attack $\epsilon_{\text{att}}=4/255$ on CIFAR-10. Table~\ref{tab:aft} reports the results. NAC provides modest additional gains on both AFT backbones ($+2.09$pp on TeCoA, $+2.44$pp on FARE), confirming that the counterattack paradigm remains beneficial even when the underlying model has been adversarially hardened. The relatively low absolute robust accuracy at this attack strength suggests that the primary bottleneck shifts from optimizer convergence to the fundamental capacity of the AFT-trained embeddings, which have been optimized under a different training objective.

\begin{table}[ht]\centering\small
\caption{AFT superposition on CIFAR-10 under PGD $\epsilon_{\text{att}}=4/255$. NAC provides complementary gains on top of adversarially fine-tuned models, confirming orthogonality between weight-space and inference-time defenses.}\label{tab:aft}
\begin{tabular}{@{}lrrrr@{}}
\toprule Base & Clean & TTC & NAC & Gain (pp) \\
\midrule TeCoA~\cite{mao2023understanding} & 63.85 & 2.75 & \textbf{4.84} & +2.09 \\
FARE~\cite{schlarmann2024robust} & 73.67 & 0.86 & \textbf{3.30} & +2.44 \\
\bottomrule
\end{tabular}
\end{table}

\subsection{Extended Iteration: NAC-4 vs.\ TTC-4}

Table~\ref{tab:nac4} provides a direct comparison of 4-step NAC and 4-step TTC across all six datasets, complementing the N-step scaling analysis in Table~\ref{tab:nstep}. Both methods benefit from additional iterations, but NAC extracts substantially greater value from each step. On CIFAR-10, NAC-4 achieves $37.05\%$ ($+11.68$pp over TTC-4); on STL-10, NAC-4 reaches $66.74\%$ ($+18.45$pp over TTC-4). These margins are considerably larger than those at $K=2$, consistent with the $O(1/K^2)$ convergence behavior of NAG---each additional iteration provides proportionally more benefit under NAC than under PGD's $O(1/K)$ regime.

\begin{table}[ht]\centering\small
\caption{Four-step counterattack comparison under PGD $\epsilon_{\text{att}}=4/255$. NAC-4 consistently outperforms TTC-4 by large margins, with the gap widening relative to the 2-step setting, consistent with the quadratic acceleration of NAG.}\label{tab:nac4}
\begin{tabular}{@{}lrrr@{}}
\toprule Dataset & TTC-4 & NAC-4 & Gain (pp) \\
\midrule CIFAR-10 & 25.37 & \textbf{37.05} & +11.68 \\
STL-10 & 48.29 & \textbf{66.74} & +18.45 \\
\bottomrule
\end{tabular}
\end{table}

\subsection{Comparison with DOC}

DOC~\cite{jiang2026diversifying} (AAAI 2026 Oral) is the strongest published extension of TTC, introducing four supplementary mechanisms---orthogonal gradient directions, directional sensitivity analysis, learnable tau gating, and gradient normalization---that collectively broaden the counterattack search space. To ensure a rigorous comparison, we evaluate all three methods (TTC, NAC, DOC) under two configurations: the NAC paper's default setting (K=2, $\alpha=1/255$) and the DOC paper's default setting (K=4, $\alpha=3/255$, learnable\_tau=0.155, temperature=75.0). In both configurations, the attack (PGD-10, $\epsilon_{\text{att}}=4/255$), model (ViT-B/32), and dataset (CIFAR-10) are held constant, isolating the counterattack algorithm as the sole variable. Table~\ref{tab:doc} reports the results.

\begin{table}[ht]\centering\small
\caption{Fair comparison of TTC, NAC, and DOC under identical conditions on CIFAR-10. In the NAC default configuration (K=2, $\alpha=1/255$), NAC substantially outperforms both TTC and DOC. In the DOC default configuration (K=4, $\alpha=3/255$), NAC achieves parity with DOC while requiring none of DOC's auxiliary mechanisms.}\label{tab:doc}
\begin{tabular}{@{}lrrr@{}}
\toprule
Method & Config A (K=2, $\alpha{=}1/255$, seed=0) & Config B (K=4, $\alpha{=}3/255$, seed=1) \\
\midrule
TTC & 7.07 & 34.64 \\
NAC & \textbf{16.47} & 38.72 \\
DOC & 4.44 & \textbf{38.81} \\
\bottomrule
\end{tabular}
\end{table}

Two findings emerge. First, in the small-iteration regime (Config~A), NAC improves upon TTC by $+9.40$pp while DOC degrades by $-$2.63pp relative to TTC. DOC's auxiliary mechanisms, designed to enhance exploration diversity, appear to introduce noise that outweighs their benefit when the iteration budget is tight. Second, with additional iterations and a larger step size (Config~B), all methods benefit substantially---TTC reaches $34.64\%$ and both NAC and DOC converge to approximately $38.7$--$38.8\%$, a difference within the $\pm 0.2$pp reproducibility margin. The practical implication is noteworthy: NAC achieves performance on par with the state-of-the-art DOC using only a single principled modification (look-ahead gradient evaluation), whereas DOC requires four additional components totaling over 30 lines of specialized logic. This result reinforces the central thesis of our work: optimizer design constitutes an underappreciated yet highly leveraged axis for improving test-time defense.

\subsection{Reproducibility}

All experiments are conducted with three random seeds ($0, 1, 2$). Across the main experimental configurations, NAC achieves $16.68\% \pm 0.19\%$ robust accuracy on CIFAR-10 under PGD $\epsilon_{\text{att}}=4/255$, compared to TTC's $7.04\% \pm 0.08\%$. The maximum observed standard deviation across all experiments is below $0.2$ percentage points, indicating stable convergence and high reproducibility.

% ============================================================
\section{Discussion}\label{sec:discussion}

\subsection{Why NAC Works: Local Smoothness and Look-Ahead}

The effectiveness of NAC rests on an empirical regularity of CLIP's ViT embedding space: under the small perturbation budget $\epsilon = 4/255$, the counterattack objective $g(\delta) = \|f_I(x_{\text{adv}} + \delta) - f_I(x)\|^2$ is locally well-behaved. In this regime, the gradient at the look-ahead position $\delta_t + \mu v_t$ provides a substantially more reliable estimate of the gradient at the next iterate than does the gradient at the current position $\delta_t$. This advantage is particularly pronounced when the optimization trajectory exhibits consistent directional structure---precisely the condition that holds when counterattacking toward a fixed clean embedding target.

Figure~\ref{fig:pca} provides qualitative evidence through PCA visualization of CLIP embeddings on CIFAR-10. Clean image embeddings (Figure~\ref{fig:pca}a) form well-separated, compact clusters corresponding to semantic classes. Under adversarial attack (Figure~\ref{fig:pca}b), these clusters collapse and scatter, with embeddings pushed toward the origin and class separation largely destroyed. TTC (Figure~\ref{fig:pca}c) partially restores class separation, but residual mixing between classes remains visible. NAC (Figure~\ref{fig:pca}d) produces markedly more compact and well-separated clusters, consistent with the quantitative gains observed in Tables~\ref{tab:main}--\ref{tab:nac4}. The visual evidence aligns with the mechanistic account: by taking more effective gradient steps, NAC more completely reverses the adversarial perturbation's effect on the embedding geometry.

\begin{figure}[ht]\centering
\includegraphics[width=0.95\textwidth]{fig4_pca.pdf}
\caption{PCA visualization of CLIP ViT-B/32 embeddings on CIFAR-10 test set. (a) Clean images: well-separated, compact class clusters. (b) Adversarial images ($\epsilon_{\text{att}}=4/255$): collapsed and scattered embeddings with severely degraded class structure. (c) After TTC defense: partial restoration of class separation. (d) After NAC defense: more complete restoration, with clusters closer to the clean embedding distribution.}\label{fig:pca}
\end{figure}

\subsection{Generality of the Optimizer Insight}

An important implication of our work extends beyond the specific NAC algorithm. DBD~\cite{liu2026directional}, CSR~\cite{nie2025contrastive}, AGC~\cite{li2026agc}, and indeed any gradient-based test-time defense that currently employs standard PGD updates can potentially benefit from replacing the underlying optimizer with Nesterov accelerated gradient. NAC's Nesterov formulation is a drop-in acceleration module: it modifies only the gradient evaluation position and velocity update, leaving the defense-specific loss function and projection step unchanged. This generality carries practical significance: many optimizers from the convex and stochastic programming literature, including accelerated and variance-reduced methods, have yet to be evaluated in the test-time defense setting.

\subsection{Why Small Iteration Budgets Amplify the Optimizer Effect}

An intriguing aspect of our findings is that the optimizer choice matters considerably more for test-time defense than for standard training. In deep learning training, where optimizers apply hundreds or thousands of iterations, the practical distinction between $O(1/K)$ and $O(1/K^2)$ convergence is often obscured by stochasticity, learning rate schedules, and the diminishing returns of late-stage optimization~\cite{kingma2015adam}. In contrast, test-time defense operates at $K=2$--$4$, a regime where each individual gradient step carries disproportionate weight and the constant factors in convergence bounds dominate. To our knowledge, this regime has not been previously studied from an optimization perspective in the context of adversarial defense. Our results suggest that small-iteration optimization constitutes a distinct design space with its own principles, meriting dedicated investigation.

\subsection{Relationship to DOC and Other Extensions}

DOC~\cite{jiang2026diversifying} (AAAI 2026 Oral) extends TTC through a different dimension: exploration diversity. DOC introduces four additional mechanisms---orthogonal gradient directions, directional sensitivity analysis, sigmoid-based learnable tau gating, and gradient normalization---that broaden the search space of the counterattack. NAC, in contrast, extends TTC through convergence rate: it accelerates descent within the existing search space.

These two extensions are fundamentally complementary, yet our empirical comparison (Section~\ref{sec:experiments}, Table~\ref{tab:doc}) reveals an instructive asymmetry. Under the NAC paper's default configuration (K=2, $\alpha=1/255$), NAC improves TTC by $+9.40$pp while DOC underperforms TTC by $-$2.63pp---suggesting that DOC's exploration mechanisms introduce noise that outweighs their benefit when the iteration budget is severely constrained. Under DOC's own optimal configuration (K=4, $\alpha=3/255$), NAC achieves $38.72\%$, statistically equivalent to DOC's $38.81\%$---yet NAC accomplishes this with a single modification (look-ahead gradient evaluation) versus DOC's four auxiliary components. This finding reinforces the efficiency of the optimizer-centric approach: a principled change in where gradients are evaluated can match the performance of substantially more complex exploration strategies.

A natural direction for future work is to combine both perspectives, equipping DOC's diversified exploration with NAC's accelerated convergence. Given NAC's demonstrated ability to match DOC's performance independently, a combined approach may yield gains beyond what either method achieves alone, particularly at intermediate iteration budgets.

\subsection{Limitations and Future Work}

Several limitations of the current work warrant discussion.

\textbf{Cross-architecture transfer.} The counterattack hyperparameters ($\epsilon_{\text{ttc}}$, $\tau$, $\beta$, step size) were tuned on ViT-B/32, and the near-zero robust accuracy on RN50 indicates that these parameters do not transfer to ResNet-based architectures. We note that this limitation is shared by prior counterattack methods: neither TTC~\cite{xing2025clip} nor DOC~\cite{jiang2026diversifying} reports cross-architecture results on RN50, and the attention-pooled embedding mechanism of ResNet-based CLIP differs substantially from ViT's patch-based features. We leave systematic hyperparameter adaptation for non-ViT backbones to future work.

A further consideration is the mismatch between the adversary's objective and the counterattack objective. Adversarial attacks minimize cross-entropy classification loss, whereas the counterattack minimizes $L_2$ embedding distance between adversarial and clean representations. Since cross-entropy operates on the logit space while the counterattack operates in the embedding space, the look-ahead gradient may correct embedding deviations without fully neutralizing the adversarial signal in the classification head. Whether a joint objective spanning both spaces can further improve robust accuracy remains an open question. On the evaluation side, our study uses ImageNet-100 rather than full ImageNet-1K---a trade-off adopted by prior CLIP robustness work~\cite{mao2023understanding,xing2025clip} that preserves class diversity while keeping extensive ablation studies tractable. Extending NAC to more recent VLMs such as BLIP-2 and LLaVA, and to tasks beyond classification, is a natural next step.

\textbf{Theoretical analysis.} As noted in Section~\ref{sec:method}, NAG's optimal $O(1/K^2)$ rate is proved for convex smooth objectives~\cite{nesterov1983method}, and a rigorous non-convex analysis for neural embedding spaces is not yet available. Ghadimi and Lan~\cite{ghadimi2016accelerated} provide guarantees for randomized NAG in non-convex stochastic programming, but their scheme requires averaging over multiple runs and variance reduction techniques incompatible with single-instance, small-$K$ test-time defense. Providing a non-asymptotic analysis of deterministic NAG under the local smoothness conditions that empirically characterize the counterattack setting is a challenging open direction.

% ============================================================
\section{Conclusion}\label{sec:conclusion}

In this paper, we have presented Nesterov Accelerated Counterattack (NAC), a training-free test-time defense for vision-language models that replaces the standard PGD optimizer with Nesterov accelerated gradient. The modification is minimal in implementation---requiring only a change in where gradients are evaluated---yet is grounded in the well-established accelerated convergence properties of NAG. NAC introduces zero additional computational cost: the forward and backward passes per iteration are identical in number and complexity to those of the standard TTC framework.

Across six benchmark datasets spanning general object recognition, fine-grained classification, and texture analysis, NAC consistently outperforms TTC under PGD attacks at three perturbation strengths ($\epsilon_{\text{att}} = 1/255, 2/255, 4/255$) and under the AutoAttack ensemble. On CIFAR-10 at $\epsilon_{\text{att}}=4/255$, NAC with two counterattack steps improves robust accuracy from $6.98\%$ to $16.61\%$ ($+9.63$ percentage points); with four steps, NAC reaches $37.05\%$ compared to TTC's $25.37\%$. Controlled ablation experiments isolate Nesterov's look-ahead mechanism from standard momentum accumulation, establishing that the look-ahead gradient evaluation---not momentum alone---accounts for the observed gains. NAC further generalizes across ViT architectures and yields complementary improvements when combined with adversarially fine-tuned models.

Beyond the specific NAC algorithm, this work identifies optimizer selection as a previously overlooked design dimension in test-time adversarial defense. The conventional practice of defaulting to standard PGD for counterattack optimization leaves substantial performance on the table. Many optimizers from the convex and stochastic programming literature---including adaptive restart schemes, variance-reduced methods, and quasi-Newton approximations---remain unexplored in the context of test-time defense. We hope that our findings will encourage the community to look beyond the specific choice of defense objective and consider the optimizer itself as a first-class design element, opening a new axis along which test-time robustness can be systematically improved.

% ============================================================
% Backmatter
\backmatter

\bmhead{Supplementary information}

The source code for reproducing all experimental results is available at \url{https://github.com/1011245276/NAC}.

\bmhead{Acknowledgements}

Supported by [funding information to be added].

\section*{Declarations}

\begin{itemize}
\item \textbf{Funding:} [To be added.]
\item \textbf{Conflict of interest:} The authors declare that they have no conflict of interest.
\item \textbf{Ethics approval and consent to participate:} Not applicable.
\item \textbf{Consent for publication:} Not applicable.
\item \textbf{Data availability:} All datasets used in this study are publicly available (CIFAR-10/100, STL-10, Flowers-102, DTD, ImageNet).
\item \textbf{Materials availability:} Not applicable.
\item \textbf{Code availability:} The source code is available at \url{https://github.com/1011245276/NAC}.
\item \textbf{Author contribution:} Zhihao Zhang: conceptualization, methodology, software, validation, writing---original draft. Yazhi Liu: supervision, writing---review and editing, funding acquisition.
\end{itemize}

% ============================================================
% Bibliography
\bibliography{nac_references}

\end{document}
